% =====================================================================
\section{Results on LHCb 2025}\label{sec:results}
% =====================================================================

This section reports the kernel fit on the LHCb 2025 dataset
\citep{LHCb2025BKstmumuComp,LHCb2025BKstmumuHEPData}, both on the
single-observable $P_{5}'$ and on the four-observable joint fit
$(P_{5}', P_{4}', P_{1}, P_{2})$. We report the headline non-linear
result first; the linearised Mode-B numbers are then quoted as a
methodology diagnostic and the discrepancy is discussed.

\subsection{Joint fit, four observables (non-linear, headline)}\label{subsec:joint_fit}

We fit the four optimised $P_{i}$ observables jointly with the full
LHCb $32{\times}32$ correlation matrix:
$(P_{5}', P_{4}', P_{1}, P_{2})$, 8 bins each, 32 data points total.
Fits are evaluated using \texttt{flavio.np\_prediction} directly
(no Taylor expansion), with one fitted parameter per model
($k=1$). Results in Table~\ref{tab:joint_fit_nl}.

\begin{table}[h]
\centering
\small
\setlength{\tabcolsep}{4pt}
\begin{tabular}{lcrrrl}
\toprule
Model & $k$ & $\chi^{2}$ & AIC & $\Delta\mathrm{AIC}$ & best fit \\
\midrule
SM                              & 0 & 160.6 & 160.6 & $+119.7$ & -- \\
\textbf{FREE\_C9} (non-linear)  & 1 &  \textbf{40.89} &  \textbf{42.89} & $\boldsymbol{0.00}$ & $\Delta C_{9}=-1.00$ \\
VFD\_GREEN\_600CELL (non-linear)& 1 &  41.98 &  43.98 & $+1.09$ & $A=+1.14$, $\Delta C_{9}^{\mathrm{eff}}=-0.86$ \\
\bottomrule
\end{tabular}
\caption{Headline non-linear joint fit on the four LHCb 2025 angular
observables (32 data points, full covariance). At equal parameter cost
the constant Wilson-coefficient shift $\mathrm{FREE\_C9}$ has lower
AIC by $1.09$ units, a marginal preference (Akaike weight ratio
$\approx 1.7$). The geometry-derived kernel produces a comparable
$\chi^{2}$ but does not beat $\mathrm{FREE\_C9}$ in this analysis.
Reproducible from \texttt{scripts/wo016c\_nonlinear\_refit.py}.}
\label{tab:joint_fit_nl}
\end{table}

\begin{figure}[h]
\centering
\includegraphics[width=0.95\textwidth]{fig_F2_bin_pulls.pdf}
\caption{Per-bin pulls $(O_{\mathrm{obs}} - O_{\mathrm{pred}})/\sigma$
on the LHCb 2025 four-observable joint fit, evaluated under the
non-linear $\mathrm{FREE\_C9}$ best fit ($\Delta C_{9}=-1.00$) and
the non-linear VFD best fit ($A=+1.14$). Dashed lines mark the
$\pm 1\sigma$ bands. Both models leave residuals of $\lesssim 2\sigma$
in most bins; neither shows a systematic spatial structure that the
other captures. Reproducible from
\texttt{scripts/wo017\_paper\_figures.py}.}
\label{fig:F2_bin_pulls}
\end{figure}

\subsection{Mode-B linearised fit (diagnostic only)}

The earlier project iteration used the Mode-B linearised response of
Eq.~\eqref{eq:modeB}, fitting $\Delta C_{9}$ via the central-difference
slope at $\delta=\pm 0.5$ around $C_{9}^{\mathrm{SM}}$. On the same
joint fit this gives Table~\ref{tab:joint_fit_linear}.

\begin{table}[h]
\centering
\small
\setlength{\tabcolsep}{4pt}
\begin{tabular}{lcrrrl}
\toprule
Model & $k$ & $\chi^{2}$ & AIC & $\Delta\mathrm{AIC}$ & best fit \\
\midrule
SM                          & 0 & 160.6 & 160.6 & $+119.3$ & -- \\
FREE\_C9 (Mode B)           & 1 &  39.30 &  41.30 & $0.00$  & $\Delta C_{9}=-1.34$ \\
VFD\_GREEN\_600CELL (Mode B)& 1 &  37.63 &  39.63 & $-1.67$ & $A=+1.59$, $\Delta C_{9}^{\mathrm{eff}}=-1.37$ \\
\bottomrule
\end{tabular}
\caption{Linearised Mode-B fit on the same data. Under the linear
model $\mathrm{VFD\_GREEN\_600CELL}$ has lower AIC by $1.67$ units.
This number does not survive non-linear evaluation: see
Table~\ref{tab:joint_fit_modes}.}
\label{tab:joint_fit_linear}
\end{table}

\subsection{Linearised vs non-linear: drift diagnostic}

\begin{table}[h]
\centering
\small
\begin{tabular}{lrrr}
\toprule
quantity & linearised (Mode B) & non-linear refit & drift \\
\midrule
$\chi^{2}$ FREE\_C9                & 39.30 & 40.89 & $+1.59$ \\
$\chi^{2}$ VFD\_GREEN\_600CELL     & 37.63 & 41.98 & $+4.35$ \\
$\Delta\mathrm{AIC}$ (VFD vs FREE) & $-1.67$ & $+1.09$ & $\mathbf{+2.77}$ \\
$\Delta C_{9}^{\mathrm{FREE}}$ best fit & $-1.34$ & $-1.00$ & $+0.34$ \\
$A$ best fit (VFD)                 & $+1.59$ & $+1.14$ & $-0.45$ \\
max per-bin linearisation residual & --  & $4.27\sigma$ & -- \\
mean per-bin linearisation residual & -- & $0.71\sigma$ & -- \\
\bottomrule
\end{tabular}
\caption{Comparison of the same fit under Mode-B linearisation vs
direct \texttt{flavio.np\_prediction}. The drift in
$\Delta\mathrm{AIC}$ is $2.77$ units, larger than the Mode-B
preference itself: the headline flips sign under non-linear
evaluation. Per-bin linearisation residuals reach $4\sigma$,
confirming that $|\Delta C_{9}|\sim 1.4$ is well outside the
linear regime. Reproducible from
\texttt{scripts/wo016c\_nonlinear\_refit.py}
(reports/wo016c\_nonlinear\_refit.md).}
\label{tab:joint_fit_modes}
\end{table}

\subsection{Single-observable fit on $P_{5}'$ (linearised, retained as a check)}

For comparability with the project's earlier reports we retain the
single-observable $P_{5}'$ fit under the linearised Mode-B response.
On the canonical 8-bin grid with the full $8{\times}8$ covariance:
SM $\chi^{2}=17.3$, FREE\_C9 $\chi^{2}=14.91$ ($\Delta C_{9}=-1.38$),
KAPPA\_EXP (continuum, Layer 1) $\chi^{2}=7.91$ at $A=+4.49$, and
VFD\_GREEN\_600CELL $\chi^{2}=13.55$ at $A=+1.68$. The continuum
exponential dominates on this single observable
($\Delta\mathrm{AIC}=-7.00$ vs FREE\_C9), reflecting that $P_{5}'$ is
the sharpest $C_{9}$-sensitive observable in this dataset and prefers
the most-peaked kernel; this preference is reduced when the four
$P$-basis observables are fitted jointly. Single-observable fits are
not run in non-linear mode because the conclusion (the kernel is
consistent with the shape of $P_{5}'$) is preserved across
both modes.

\subsection{Leave-one-observable-out stability (linearised, retained)}

Refitting the shared kernel under Mode B after removing one of the
four observables gives Table~\ref{tab:loo_obs}.

\begin{table}[h]
\centering
\small
\begin{tabular}{lcccc}
\toprule
dropped & $n$ remaining & $A$ (Mode B) & $\Delta C_{9}^{\mathrm{eff}}$ mean & $\chi^{2}$ \\
\midrule
$P_{1}$  & 24 & 1.558 & $-1.336$ & 29.23 \\
$P_{2}$  & 24 & 1.651 & $-1.416$ & 27.70 \\
$P_{4}'$ & 24 & 1.620 & $-1.390$ & 33.46 \\
$P_{5}'$ & 24 & 1.581 & $-1.356$ & 26.86 \\
\midrule
all (reference, Mode B)  & 32 & 1.594 & $-1.367$ & 37.63 \\
\bottomrule
\end{tabular}
\caption{Leave-one-observable-out fits of the shared kernel under
Mode B (linearised). $A$ varies within $\pm 3\,\%$ across the four
folds, $\Delta C_{9}^{\mathrm{eff}}$ stays in $[-1.42, -1.34]$, sign
always negative. The non-linear LOO produces qualitatively the same
sign-stable behaviour at lower magnitudes
($A\in[1.05, 1.20]$); we retain the linearised values here for
continuity with the project's earlier multi-observable joint-fit
report.}
\label{tab:loo_obs}
\end{table}

\subsection{Acceptance gates}

The acceptance criteria pre-committed for this multi-observable
joint-fit test were stated for the linearised analysis. Under the
headline non-linear analysis they read as follows:
\begin{enumerate}[label=(\roman*)]
\item $\mathrm{VFD\_GREEN\_600CELL}\ \mathrm{AIC} \leq \mathrm{FREE\_C9}\ \mathrm{AIC}$:
  \textsc{fail} (non-linear $\Delta\mathrm{AIC}=+1.09$, marginal in
  the opposite direction).
\item All major angular observables prefer the same sign of effective
  $\Delta C_{9}$ in the shared-amplitude fit:
  \textsc{pass} ($A>0$ on every fold; $\Delta C_{9}^{\mathrm{eff}}<0$
  on every bin in the headline non-linear fit).
\item Leave-one-observable-out remains stable in sign and bounded in
  magnitude: \textsc{pass} ($A$ within $\pm 3\,\%$ across folds, sign
  always positive).
\end{enumerate}

Acceptance criterion (i) does not pass under the headline non-linear
analysis on this single dataset; the kernel is statistically
indistinguishable from a constant Wilson-coefficient shift on AIC.
Sign-stability and amplitude-bounded behaviour survive both modes.
The next two sections establish whether the kernel transfers in a
sign-uniform way to other datasets (\S\ref{sec:cross-dataset}) and
to other channels (\S\ref{sec:cross-channel}), under the full
non-linear analysis.
